\section{Virasoro algebra}\label{virasoro-algebra}

\subsection{Asymptotic Killing vectors}\label{asymptotic-killing-vectors}

Starting from the exact Killing fields \(\xi _{a}\), \(\bar{\xi}_{a}\) and the large-diffeomorphism parameters \(\zeta _{-,0}^{\mu}\) and \(\zeta _{+,0}^{\mu}\) found in Section 2, we define the first non-isometric asymptotic generators by

\begin{align}
\xi _{-2} & =\frac{3}{2i\sqrt{ 3G }}\zeta _{-,0}, & \xi _{2} & =\xi _{-2}^{*} \\
\bar{\xi}_{-2} & =\frac{3}{2i\sqrt{ 3G }}\zeta _{+,0}, & \bar{\xi}_{2} & =\bar{\xi}_{-2}^{*}
\end{align}

Higher asymptotic generators are defined recursively as

\begin{align}
\xi _{n+1} & =\frac{1}{i(n-1)}\mathcal{L}_{\xi_{1}}\xi _{n}, & \xi _{-(n+1)} & =-\frac{1}{i(n-1)}\mathcal{L}_{\xi _{-1}}\xi _{-n} \\
\bar{\xi}_{n+1} & =\frac{1}{i(n-1)}\mathcal{L}_{\bar{\xi}_{1}}\bar{\xi}_{n}, & \bar{\xi}_{-(n+1)} & =-\frac{1}{i(n-1)}\mathcal{L}_{\bar{\xi}_{-1}}\bar{\xi}_{-n}, & n\geqslant 2
\end{align}

These asymptotic generators obey two copies of the Witt algebra asymptotically,

\begin{align}
[\xi _{n},\xi _{m}] & =-i(n-m)\xi _{n+m}+\mathcal{O}(r^{-4}) \\
[\bar{\xi}_{n},\bar{\xi}_{m}] & =-i(n-m)\bar{\xi}_{n+m}+\mathcal{O}(r^{-4}) \\
[\xi _{n},\bar{\xi}_{m}] & =0,\quad n,m\in \mathbb{Z}
\end{align}

\subsection{Noether theorem for (asymptotic) Killing symmetries}\label{noether-theorem-for-asymptotic-killing-symmetries}

We now discuss the Noether charges associated with the asymptotic Killing symmetries. Such a symmetry is represented on configuration space by the vector field

\begin{align}
X_{\xi} & =\int \mathrm{d}^{3}x\mathcal{L}_{\xi}g_{\mu \nu} \frac{\delta}{\delta g_{\mu \nu}} \\
 & =\int \mathrm{d}^{3}x\left(\mathcal{L}_{\xi}g^{(0)}_{\mu \nu} +\mathcal{L}_{\xi}h_{\mu \nu}+\mathcal{O}(h^{2})\right) \frac{\delta}{\delta h_{\mu \nu}}
\end{align}

Here \(\xi ^{\mu}\) is an asymptotic Killing vector: \(\nabla _{\mu}\xi _{\nu}+\nabla _{\nu}\xi _{\mu}\) obeys the Brown-Henneaux falloffs. The vector field itself has the asymptotic behavior

\begin{align}
\xi ^{t} & =\mathcal{O}(r^{0}), & \xi ^{r} & =\mathcal{O}(r), & \xi ^{\phi} & =\mathcal{O}(r^{0})
\end{align}

and preserves the Brown-Henneaux falloffs of the perturbation:

\begin{align}
X_{\xi}\cdot \delta h_{tt} & =\mathcal{O}(r^{0}), & X_{\xi}\cdot \delta h_{tr} & =\mathcal{O}(r^{-3}), & X_{\xi}\cdot \delta h_{t\phi} & =\mathcal{O}(r^{0}) \\
 & & X_{\xi}\cdot \delta h_{rr} & =\mathcal{O}(r^{-4}), & X_{\xi}\cdot \delta h_{r\phi} & =\mathcal{O}(r^{-3}) \\
 &  &  &  & X_{\xi}\cdot \delta h_{\phi \phi} & =\mathcal{O}(r^{0})
\end{align}

Acting with \(X_{\xi}\) on the bulk action gives

\begin{align}
X_{\xi}\cdot \delta S_{\text{bulk}} & =\alpha _{\xi}|_{\Sigma _{f}}-\alpha _{\xi}|_{\Sigma _{i}} \\
\alpha _{\xi} & =-\frac{1}{16\pi G}\int _{\Sigma}\mathrm{d}^{2}x\sqrt{ \sigma ^{(0)} }\tau ^{(0)}_{\mu}\xi ^{\mu}\left(\nabla ^{(0)}_{\rho}\nabla ^{(0)}_{\sigma}h^{\rho \sigma}-\nabla ^{(0)2}h+\mathcal{O}(h^{2})\right)
\end{align}

The corresponding Noether charge is

\begin{align}
H_{\xi} & =X_{\xi}\cdot \theta-\alpha _{\xi} \\
 & \approx \int _{\Sigma}\mathrm{d}^{2}x\sqrt{ \sigma ^{(0)} }\tau ^{(0)} _{\mu}\nabla ^{(0)} _{\nu}Q_{\xi}^{\nu \mu} \\
 & =\int _{\partial \Sigma}\mathrm{d}x\sqrt{ h^{(0)} }\tau _{\mu}^{(0)}n^{(0)}_{\nu}Q_{\xi}^{\mu \nu}
\end{align}

Here \(h^{(0)}_{ij}\) is the metric induced on \(\partial \Sigma=\Sigma \cap \Gamma\),

\begin{align}
\mathrm{d}s^{2}_{\partial \Sigma} & =h^{(0)}_{ij}\mathrm{d}x^{i}\mathrm{d}x^{j} \\
 & =r^{2}\mathrm{d}\phi ^{2}, \quad r\to \infty
\end{align}

and \(n^{(0)\mu}\) is the outward-pointing unit normal vector orthogonal to the spatial boundary \(\Gamma\):

\begin{align}
n^{(0)\mu} & =\sqrt{ 1+r^{2} }\delta ^{\mu}_{r}
\end{align}

The antisymmetric potential \(Q_{\xi}^{\mu \nu}\) is

\begin{align}
Q_{\xi}^{\mu \nu} & =\frac{1}{16\pi G}\left(A^{(1)\mu \nu}_{\xi}+\frac{1}{2}hA_{\xi}^{(0)\mu \nu}\right)+\mathcal{O}(h^{2}) \\
A_{\xi}^{(0)\mu \nu} & =\nabla ^{(0)\mu}\xi ^{\nu}-\nabla ^{(0)\nu}\xi ^{\mu} \\
A_{\xi}^{(1)\mu \nu} & =-h^{\mu \rho}\nabla ^{(0)}_{\rho}\xi ^{\nu}+h^{\nu \rho}\nabla ^{(0)}_{\rho}\xi ^{\mu}+\xi ^{\rho}(\nabla ^{(0)\mu}h^{\nu}_{~\rho}-\nabla ^{(0)\nu}h^{\mu}_{~\rho})
\end{align}

Higher-order contributions will not affect the central term computed below. To compute the Poisson brackets between Noether charges, we consider the variation of \(H_{\xi}\):

\begin{align}
\delta H_{\xi} & \approx \int _{\partial \Sigma}\mathrm{d}x\sqrt{ h^{(0)} }\tau ^{(0)} _{\mu}n^{(0)}_{\nu}\left(\delta Q_{\xi}^{\mu \nu}-2\xi ^{[\mu}\theta ^{\nu]}\right) \\
 & =\int _{\partial \Sigma}\mathrm{d}x\sqrt{ h^{(0)} }\tau _{\mu}^{(0)}n_{\nu}^{(0)}k _{\xi}^{\mu \nu}
\end{align}

where \(\xi ^{[\mu}\theta ^{\nu]}=\frac{1}{2}(\xi ^{\mu}\theta ^{\nu}-\xi ^{\nu}\theta ^{\mu})\), and

\begin{align}
k _{\xi}^{\mu \nu} & =\delta Q_{\xi}^{\mu \nu}-2\xi ^{[\mu}\theta ^{\nu]} \\
 & =\frac{1}{16\pi G}\left[ \xi ^{\rho}(\nabla ^{(0)\mu}\delta h^{\nu}_{~\rho}-\nabla ^{(0)\nu}\delta h^{\mu}_{~\rho})+\frac{1}{2}\delta h\left(\nabla ^{(0)\mu}\xi ^{\nu}-\nabla ^{(0)\nu}\xi ^{\mu}\right) \right. \\
 & -\delta h^{\mu \rho}\nabla ^{(0)}_{\rho}\xi ^{\nu}+\delta h^{\nu \rho}\nabla ^{(0)}_{\rho}\xi ^{\mu}+\xi ^{\mu}(\nabla ^{(0)}_{\rho}\delta h^{\nu \rho}-\nabla ^{(0)\nu}\delta h) \\
 & \left.-\xi ^{\nu}(\nabla ^{(0)}_{\rho}\delta h^{\mu \rho}-\nabla ^{(0)\mu}\delta h)\right]+\mathcal{O}(h\delta h)
\end{align}

Again, higher-order terms are dropped. The Poisson bracket between two Noether charges \(H_{\xi}\) and \(H_{\zeta}\) is computed from the covariant symplectic form as

\begin{align}
\left\{H_{\xi},H_{\zeta}\right\} & =X_{\xi}\cdot X_{\zeta}\cdot \Omega \\
 & \approx X_{\zeta}\cdot \delta H_{\xi} \\
 & =\int _{\partial \Sigma}\mathrm{d}x\sqrt{ h^{(0)} }\tau ^{(0)}_{\mu}n^{(0)}_{\nu}(X_{\zeta}\cdot k _{\xi}^{\mu \nu})
\end{align}

The weak equality means that the linearized equations of motion have been imposed.

\subsection{Virasoro algebra}\label{virasoro-algebra-1}

We now apply this charge formula to the asymptotic Killing vectors constructed above. Define

\begin{align}
X_{\xi _{n}} & =\int \mathrm{d}^{3}x\left(\mathcal{L}_{\xi _{n}}g^{(0)}_{\mu \nu}+\mathcal{L}_{\xi _{n}}h_{\mu \nu}+\mathcal{O}(h^{2})\right) \frac{\delta}{\delta h_{\mu \nu}} \\
X_{\bar{\xi}_{n}} & =\int \mathrm{d}^{3}x\left(\mathcal{L}_{\bar{\xi}_{n}}g^{(0)}_{\mu \nu}+\mathcal{L}_{\bar{\xi}_{n}}h_{\mu \nu}+\mathcal{O}(h^{2})\right) \frac{\delta}{\delta h_{\mu \nu}}
\end{align}

and let \(H_{n}\) and \(\bar{H}_{n}\) be the corresponding Noether charges. On the TT representatives used in Section 2, the charge and its variation reduce to

\begin{align}
H_{\xi} & \approx \frac{1}{16\pi G}\int _{\partial \Sigma}\mathrm{d}x\sqrt{ h^{(0)} }\tau ^{(0)}_{\mu}n^{(0)}_{\nu}\left(-h^{\mu \rho}\nabla ^{(0)}_{\rho}\xi ^{\nu}+h^{\nu \rho}\nabla ^{(0)}_{\rho}\xi ^{\mu}+\xi ^{\rho}\left(\nabla ^{(0)\mu}h^{\nu}_{~\rho}-\nabla ^{(0)\nu}h^{\mu}_{~\rho}\right)+\mathcal{O}(h^{2})\right) \\
\delta H_{\xi} & \approx \frac{1}{16\pi G}\int _{\partial \Sigma}\mathrm{d}x\sqrt{ h^{(0)} }\tau _{\mu}^{(0)}n_{\nu}^{(0)}\left(-\delta h^{\mu \rho}\nabla ^{(0)}_{\rho}\xi ^{\nu}+\delta h^{\nu \rho}\nabla ^{(0)}_{\rho}\xi ^{\mu}+\xi ^{\rho}(\nabla ^{(0)\mu}\delta h^{\nu}_{~\rho}-\nabla ^{(0)\nu}\delta h^{\mu}_{~\rho})+\mathcal{O}(h\delta h)\right)
\end{align}

We use

\begin{align}
\{H_\xi,H_\zeta\}\approx X_\zeta\cdot\delta H_\xi.
\end{align}

Substituting the asymptotic Killing vectors \(\xi _{n}\) and \(\bar{\xi}_{n}\) gives the desired Virasoro algebra

\begin{align}
\{H_n,H_m\}&=-i(n-m)H_{n+m}-\frac{i}{8G}n(n^2-1)\delta_{n+m,0},\\
\{\bar H_n,\bar H_m\}&=-i(n-m)\bar H_{n+m}-\frac{i}{8G}n(n^2-1)\delta_{n+m,0},\\
\{H_n,\bar H_m\}&=0.
\end{align}

This is the classical Virasoro algebra with central charge \(c=\frac{3}{2G}\). In this evaluation, the Witt term comes from the field-dependent part \(\mathcal{L}_{\zeta}h_{\mu\nu}\), while the field-independent central term comes from the background variation \(\mathcal{L}_{\zeta}g^{(0)}_{\mu\nu}\) in the symmetry vector field \(X_{\zeta}\).
